\section{Method}
\Missing{DESIGN ONLY: this section must be replaced by the final implemented method.}

The current candidate design question is how a predicate-conditioned interval
can explain the observed evidence without including unsupported co-presence or
discarding valid event evidence. One abstract formulation is
\begin{equation}
S(p,I\mid X)=S_{\mathrm{support}}-\lambda S_{\mathrm{extension}}
-\mu S_{\mathrm{omission}}.
\end{equation}
This equation is a design requirement, not an implemented algorithm or a novelty
claim. The completed method must define all three terms, their normalization,
training targets, inference computation, and treatment of repeated occurrences.

\paragraph{Required architecture specification.}
\Missing{Define tensor shapes, predicate conditioning, query/event assignment,
legal coordinates, interval selection, confidence, and fixed-budget inference.}

\paragraph{Training objective.}
\Missing{Write only losses used in code. Specify native initialization, frozen
parameters, matching, unknown-label handling, and train-only statistics.}

\paragraph{Ordinary boundary control.}
A start/end baseline can select legal endpoints using
\begin{equation}
\arg\max_{i<j}\left[\log p_s(i)+\log p_e(j)+
\lambda_L\log p(L=j-i\mid p)\right].
\end{equation}
Conditioned boundaries and duration priors have close precedent. Their inclusion
does not itself establish the main contribution, and true long events must retain
support throughout the duration prior.
